% Draft: Fusion methods subsection (proposed solution)
% Replaces lines 193–212 of proposed_solution-5.tex

\subsection{Fusion methods}
\label{subsec:fusion_methods_proposed}

We evaluate seven methods from three families introduced in Section \ref{sec:fusion}: SFIM \cite{liu2000smoothing} and GLP \cite{aiazzi2006mtf} from the component substitution / MRA family, CNMF \cite{yokoya2012cnmf} and CSTF \cite{li2018fusing} from the matrix and tensor factorization family, and MAPSMM \cite{hardie2004map}, HySure \cite{simoes2015hysure} and FUSE \cite{wei2015fuse} from the Bayesian / variational family.

All seven implementations are taken from the \texttt{hif-benchmarking} framework \cite{hif_review_2022} with minor modifications where relevant. Two of the seven, SFIM and GLP, are reimplemented in Python; the remaining five run in MATLAB. The evaluation results are presented in Section \ref{sec:fusion_evaluation}.

\paragraph{SFIM.}
SFIM produces the fused image by multiplying each upsampled HS band with a per-pixel ratio that carries the spatial detail of the MS image \cite{liu2000smoothing}. The ratio requires a high-resolution reference \(\mathbf{p}_k\) spectrally matched to HS band \(k\). We construct \(\mathbf{p}_k\) using the hyper-sharpening approach of Selva et al.\ \cite{selva2015hyper}: at LR, a non-negative least-squares (NNLS) regression with an intercept term is fitted from the MS bands to each HS band, and the same coefficients are then evaluated on the native-resolution MS image. The fusion follows Equation~\ref{eq:sfim_form}: the upsampled HS band is multiplied by \(\mathbf{p}_k / \widetilde{\mathbf{p}}_k\), where \(\widetilde{\mathbf{p}}_k\) is obtained by applying the same regression coefficients to the MS image bicubic-downsampled to EMIT resolution and then upsampled back to 10\,m. The ratio equals approximately 1 at LR scales and departs from 1 only at finer scales, so the multiplication injects spatial detail without altering the HS band's coarse-scale radiometry. The ratio is clamped to \([0, 2]\) to suppress noise amplification in dark pixels.

\paragraph{GLP.}
GLP replaces the multiplicative injection of SFIM with an additive high-pass residual \cite{aiazzi2006mtf}. As in SFIM, a per-band synthetic reference \(\mathbf{p}_k\) is built by ordinary least-squares regression (with intercept) from MS bands to each HS band, fitted at LR and applied at native MS resolution. A low-pass version \(\widetilde{\mathbf{p}}_k\) is obtained by filtering \(\mathbf{p}_k\) with a Gaussian kernel whose standard deviation is set so that its frequency response at the Nyquist frequency matches a target sensor MTF gain. The kernel is windowed with a Kaiser window. The fused band is
\[
  \widehat{\mathbf{z}}_k
    = \widetilde{\mathbf{y}}_{H,k}
      + g\,(\mathbf{p}_k - \widetilde{\mathbf{p}}_k),
\]
where \(\widetilde{\mathbf{y}}_{H,k}\) is the HS band upsampled by FIR interpolation and \(g = \mathrm{cov}(\widetilde{\mathbf{y}}_{H},\, \widetilde{\mathbf{p}}) \,/\, \mathrm{var}(\widetilde{\mathbf{p}})\) is a global gain computed across all bands and pixels that scales the detail amplitude to the HS radiometry.

\paragraph{CNMF.}
CNMF decomposes both the HS and MS images into a shared set of spectral endmembers and image-specific abundance maps, then reconstructs the fused image by combining the HS-derived endmembers with the MS-derived abundances \cite{yokoya2012cnmf}.

The spectral coupling between the two sensors is encoded in a spectral response matrix \(\mathbf{R} \in \mathbb{R}^{L_m \times L_h}\), where each row describes how one MS band averages across the HS spectrum. Rather than estimating \(\mathbf{R}\) from the data, we compute it analytically from the published Sentinel-2A spectral response functions and the EMIT band centres (modelled as Gaussians with 7.5\,nm FWHM):
\begin{equation}
  R_{i,j} = \frac{\int S_i(\lambda)\, E_j(\lambda)\, d\lambda}
                  {\sum_{j'} \int S_i(\lambda)\, E_{j'}(\lambda)\, d\lambda},
  \label{eq:srf_R}
\end{equation}
where \(S_i(\lambda)\) is the spectral response of MS band \(i\) and \(E_j(\lambda)\) is the response of HS band \(j\).

The number of endmembers \(M\) is chosen adaptively as the minimum of 20, the effective rank of the HS data matrix, and \(L_h - 1\), with a floor of 3. The initial endmember spectra are extracted by Vertex Component Analysis (VCA) \cite{nascimento2005vca}. The HS image is then factorised as \(\mathbf{Y}_h \approx \mathbf{W}\mathbf{H}_h\), where \(\mathbf{W} \in \mathbb{R}^{L_h \times M}\) contains the endmember spectra and \(\mathbf{H}_h\) the corresponding LR abundance maps, using multiplicative NMF updates with a sum-to-one penalty on the abundances. The MS image is factorised as \(\mathbf{Y}_m \approx (\mathbf{R}\mathbf{W})\mathbf{H}_m\), coupling the two decompositions through \(\mathbf{R}\): the MS endmember matrix is not free but is fixed to \(\mathbf{R}\mathbf{W}\), so the MS factorisation recovers abundance maps \(\mathbf{H}_m\) at the native MS resolution. The fused image is then
\[
  \widehat{\mathbf{Z}} = \mathbf{W}\,\mathbf{H}_m,
\]
combining the HS-resolution spectra with the MS-resolution abundances.

\paragraph{MAPSMM.}
MAPSMM formulates fusion as maximum a posteriori (MAP) estimation of the HR-HS image under a stochastic mixing model (SMM) prior \cite{hardie2004map}. The HS image is first projected into a PCA subspace. An SMM with several pure-material classes is fitted to the principal components, yielding class means, covariances, and per-pixel membership probabilities. The leading principal components are then enhanced by jointly using the MS image and the SMM statistics as a prior, while lower components, which carry little spatial structure, are replicated without MS guidance. The enhanced components are transformed back to the original spectral space to produce the fused image.

\paragraph{HySure.}
HySure formulates fusion as a convex optimisation problem solved by ADMM \cite{simoes2015hysure}. The fused image is represented in a low-dimensional subspace whose basis is learned from the HS data by VCA (run 20 times, keeping the highest-volume simplex). The objective enforces three terms: fidelity to the HS observations through the spatial blur kernel, fidelity to the MS observations through the spectral response matrix \(\mathbf{R}\), and spatial smoothness via total-variation regularisation on the subspace coefficients. ADMM splits the problem into closed-form sub-updates for each term. After convergence, the subspace coefficients are projected back to the full spectral space and denoised by SVD truncation.

\paragraph{FUSE.}
FUSE casts fusion as Bayesian estimation in a spectral subspace and obtains the MAP solution by solving a Sylvester matrix equation \cite{wei2015fuse}. The HS image is projected onto its leading PCA components. Under a Gaussian prior, the posterior mean reduces to a linear system whose structure (a sum of Kronecker products) admits a closed-form solution via the Sylvester equation, avoiding iterative optimisation. The spatial blur kernel is modelled as a Gaussian with FWHM equal to the scale factor. The spectral response \(\mathbf{R}\) is estimated from the data by constrained regression. The subspace solution is then mapped back to the full spectral domain.

\paragraph{CSTF.}
CSTF represents the fused image as a Tucker decomposition with three factor matrices: two spatial dictionaries learned from the MS image by K-SVD, and one spectral dictionary extracted from the HS image by VCA \cite{li2018fusing}. The core tensor and factor matrices are refined iteratively: each spatial dictionary is updated by conjugate gradient while enforcing consistency with a blurred-and-downsampled version of itself, the spectral dictionary is updated under the spectral response constraint, and the core tensor is re-estimated by sparse Tucker decomposition. After convergence, the fused image is reconstructed as the Tucker product of the three updated dictionaries and the core tensor.
